\section{Literature review}
A wide range of approaches exist to solve orbital ascent trajectory optimization problems, including direct and indirect methods using both shooting and collocation for discretization.

Despite its disadvantages indirect methods are widely used when computing space flight trajectories. The main reason is the low accuracy associated with direct methods. In practice when using direct methods the minimum functional value found can be expected to have an error of about one percent which can be a crucial part of the payload in a space flight mission. Increasing the finite-dimensional space doesn't necessarily yield better results \cite{von1992direct}.

One such indirect method is \acrfull{sgra}. Developed during the years 1968-86 by Miele \emph{et al.} \cite{miele1970sequential}, this method has been successfully applied to atmospheric flight \cite{miele1986optimal}, space flight \cite{miele1999optimaltransfers, miele2001optimaltrajectories} and flights at the interface between atmosphere and space \cite{miele1998optimal, miele1993nominal, miele2000feasibility, miele2001assessment}. Applications and extensions have been developed all around the world with a notable example being the version of this algorithm used at the NASA Johnson Space Center under the code name SeGRAm \cite{rishikof1992segram, miele2007optimal}.

\acrshort{sgra} is a sequential algorithm that alternates between gradient and restoration phases. In the gradient phase nominal functions for the states, $x(t)$, controls $u(t)$ and parameter $\pi$, satisfying all the differential equations and boundary conditions are assumed. Variations $\Delta x(t)$, $\Delta u(t)$ and $\Delta \pi$ leading to varied functions $\tilde x(t)$, $\tilde u(t)$ and $\tilde \pi$, are determined so that the value of the objective function is decreased. These variations are obtained by minimizing the first-order change of the functional subject to the linearized differential equations, the linearized boundary conditions, and a quadratic constraint on the variations of the control and the parameter. 

During this process the resulting $\tilde x(t)$, $\tilde u(t)$ and $\tilde \pi$ may not satisfy the differential equations or boundary conditions anymore. This is corrected in the restoration phase by assuming, $\tilde x(t)$, $\tilde u(t)$ and $\tilde \pi$, are the nominal functions and determining variations $\Delta \tilde x(t)$, $\Delta \tilde u(t)$ and $\Delta \tilde \pi$ that lead to varied functions $\hat x(t)$, $\hat u(t)$ and $\hat \pi$, that are consistent with the differential equations and boundary conditions.

In 2003 Miele \emph{et al.} developed an extension to \acrshort{sgra} called \acrfull{msgra}. The reason was to improve the robustness of the method and enlarge the field of applications. \acrshort{msgra} can be applied to optimal control problems involving multiple subsystems as well as discontinuities in the state and control variable at the interface between contiguous subsystems \cite{miele2003multiplep1}. \acrshort{msgra} has demonstrated in being an effective tool for optimizing the trajectories of multiple stage launch vehicles \cite{miele2003multiplep2}.

As already mentioned a mayor disadvantage of indirect methods is the high sensitivity and difficulty to select initial guesses, particularly in the case of the adjoint variables. In other to lessen this problem some applications apply the use of the homotopy method. The principle of the homotopy method is that a given problem is embedded into a family of problems parametrized by a homotopic parameter. The solution to the original problem can then be found by starting with a simpler problem and gradually adjust the homotopic parameter such that in the end a solution to the original problem is found \cite{pan2016double}.

As an example Calise \emph{et al.} \cite{calise2000further} and, Gath and Calise \cite{gath2001optimization} propose the use of a hybrid analytical/numerical approach which uses a homotopy method that starts with a nearly analytic vacuum solution that gradually introduces atmospheric effects until a converged solution is obtained. This method has the potential to be used online, since once converged, new trajectories can be computed within milliseconds using input from the guidance system due to the nearly analytic nature of the solution process \cite{gath2001optimization}. 

The issue with the homotopy method is complexity is added to the solution process.

Direct methods have also been proposed, such as those by Roh and Kim \cite{roh2002trajectory} and Jamilnia and Naghash \cite{jamilnia2012simultaneous}, both of whom solve the problem by using direct collocation to transcribe the trajectory optimization problem into \acrfull{nlp} problems. Another direct method proposed by Coverstone-Carrol and Williams \cite{coverstone1994optimal} uses differential inclusion concepts to remove explicit control dependence from the problem statement thereby reducing the number unknown parameters.

Recently newer techniques make use of heuristic optimization methods. A problem inherent to direct and indirect methods is that since these methods are based on local computations, they do not always converge towards to the global optimum. On the other hand heuristic methods are global techniques which can look for the global optimum within a larger portion of the state space. The main idea behind heuristic optimization methods is that the search is done in a stochastic manner instead of a deterministic manner \cite{rao2009survey}.

As an example Pontani \emph{et al.} \cite{PONTANI2014852, pontani2013ascent, pontani2014simple} propose a method for computing the optimal exoatmospheric trajectory of the upper stage using \acrfull{pso}. The proposed method starts by analytically deriving the necessary conditions for optimality just as with indirect methods like \acrshort{sgra}, then \acrshort{pso} is used to numerically determine the state, control and adjoint variables. This significantly lessens the initialization issues associated with indirect methods, since heuristic methods can search in a larger region of the state space for the global optimum. 

A more direct heuristic optimization approach is that taken by Zheng \emph{et al.} \cite{zheng2020ascent}. A pitch law is defined consisting of five successive sections, each for a single phase of flight with the first one starting at take-off and the last one ending at orbit insertion. A function with unknown parameters is chosen for each section of the pitch law, and then the unknowns are found using \acrfull{ga}. Since these functions can be selected based on prior knowledge, it's possible to significantly reduce the number of unknowns. In this case there where only nine.

This is a similar approach as to that taken by Pontani and Teofilatto \cite{pontani2014simple} to compute the trajectory during atmospheric flight. However in this case, the parameters where found using gradient based optimization methods instead of heuristic methods.

